\chapter{Matriz de Coerência do TCC}

A matriz de coerência ajuda a verificar se as partes do trabalho estão alinhadas. Ela pode ser usada no projeto, durante a escrita e antes da entrega.

\begin{longtable}{p{0.26\linewidth}p{0.62\linewidth}}
\toprule
Item & Pergunta de verificação \\
\midrule
\endfirsthead
\toprule
Item & Pergunta de verificação \\
\midrule
\endhead
Tema & O tema está delimitado e é compatível com o curso? \\
Problema & A pergunta de pesquisa é clara e investigável? \\
Objetivo geral & O objetivo geral responde diretamente ao problema? \\
Objetivos específicos & Os objetivos específicos conduzem ao objetivo geral? \\
Justificativa & A relevância do trabalho está demonstrada com argumentos e fontes? \\
Referencial teórico & Os conceitos utilizados ajudam a analisar o problema? \\
Metodologia & O método permite responder ao problema? \\
Dados ou corpus & Os materiais analisados são adequados e suficientes? \\
Resultados & Os resultados se relacionam aos objetivos? \\
Discussão & A análise dialoga com a literatura? \\
Conclusão & A conclusão responde ao problema e reconhece limites? \\
Referências & Todas as fontes citadas estão corretamente registradas? \\
\bottomrule
\end{longtable}

\section{Como usar}

Preencha a matriz em linguagem simples. Se algum item não puder ser respondido com segurança, o trabalho provavelmente precisa de ajuste. A matriz também pode ser apresentada ao orientador para facilitar a discussão sobre coerência do projeto.

\caixaorientacao{A matriz não é parte obrigatória do TCC final. Ela é um instrumento de planejamento e revisão.}

